\subsection{Example 1: Product Research Conversation}
\label{subsec:example-product-research}
The first example traces a complete request through the ad-serving pipeline,
from a user query about marathon footwear to an impression recorded in the
database.  Figure~\ref{fig:example-1} shows the rendered response: the
assistant naturally weaves in a \textsc{[Sponsored]} reference to the
Nike Pegasus~41 as an inline hyperlink within an otherwise helpful answer.

\begin{figure}[H]
    \centering
    \includegraphics[width=\textwidth]{images/illustrative_example_1.png}
    \caption{Chat interface response to ``What are the best running shoes for
    marathon training?''.}
    \label{fig:example-1}
\end{figure}

\subsubsection{Stage 1: Context Extraction}
\label{subsubsec:example1-context}

The orchestrator passes the raw message to the context-extraction module, which
classifies intent, estimates ad receptivity, and infers a purchase-proximity
signal.  For this query the extractor returns:

\begin{lstlisting}[style=jsonstyle, caption={Context extraction output.},
label={lst:example1-context}]
{
  "intent":          "product_research",
  "ad_receptivity":  "medium",
  "purchase_signal": 0.80,
  "topics":          ["running shoes", "marathon training"],
  "entities":        []
}
\end{lstlisting}

\noindent
A \texttt{purchase\_signal} of~0.80 boosts commercially relevant candidates in
the re-ranker.  An \texttt{ad\_receptivity} of \texttt{medium} allows the
pipeline to proceed; \texttt{none} would skip retrieval entirely.

\subsubsection{Stage 2: Vector Search and Re-ranking}
\label{subsubsec:example1-retrieval}

The topics are concatenated into a retrieval string and encoded into a dense
vector.  A nearest-neighbour search over the \texttt{pgvector} embeddings
column returns three candidates ranked by cosine similarity, with their
campaign budget fields attached for the gate check:

\begin{lstlisting}[style=jsonstyle, caption={Vector-search candidates returned
by the retrieval stage.}, label={lst:example1-retrieval}]
[
  {
    "ad_id":        "ad_nike_pegasus_41",
    "product_name": "Nike Pegasus 41",
    "similarity":   0.7918,
    "bid_cpm":      5.00,
    "daily_budget": 100.00,
    "today_spend":  12.40
  },
  {
    "ad_id":        "ad_allbirds_tree_runners",
    "product_name": "Allbirds Tree Runners",
    "similarity":   0.6961,
    "bid_cpm":      4.50,
    "daily_budget": 80.00,
    "today_spend":  18.20
  },
  {
    "ad_id":        "ad_notion_team_plan",
    "product_name": "Notion Team Plan",
    "similarity":   0.4991,
    "bid_cpm":      3.00,
    "daily_budget": 60.00,
    "today_spend":  5.10
  }
]
\end{lstlisting}

\noindent
The re-ranker combines similarity, bid CPM, and purchase signal into a
composite score.  The Nike Pegasus~41 ranks first ($\text{similarity}=0.7918$),
followed by the Allbirds Tree Runners ($0.6961$) and the Notion Team Plan
($0.4991$).  The Notion Team Plan is filtered out by the topic-relevance
threshold, as productivity software is semantically distant from the running
footwear query.  The budget gate confirms that
$12.40 + 0.005 \ll 100.00$ for the top candidate, so the Nike Pegasus~41
passes and becomes the sole injection candidate.

\subsubsection{Stage 3: Prompt Assembly}
\label{subsubsec:example1-prompt}

The winning ad is inserted into the synthesis prompt via a
\texttt{<sponsored\_context>} block carrying the structured ad fields and
its \emph{context guide}---a free-text presentation directive refined by the
adaptive optimiser (see Section~\ref{subsec:example-context-optimization}):

\begin{lstlisting}[style=jsonstyle, caption={Sponsored context block injected
into the system prompt.}, label={lst:example1-prompt}]
<sponsored_context>
  <ad id="ad_nike_pegasus_41">
    <product>Nike Pegasus 41</product>
    <description>Responsive ZoomX foam; breathable Flyknit upper;
      ideal for daily training and long-distance runs.</description>
    <tracking_url>https://adservice.example/click/ad_nike_pegasus_41
    </tracking_url>
    <context_guide>
      Integrate naturally as one option among several.  Always prefix
      with [Sponsored].  Connect ZoomX cushioning and breathable upper
      to the user's stated training goals.
    </context_guide>
  </ad>
</sponsored_context>
\end{lstlisting}

\subsubsection{Stage 4: Response Generation and Impression Logging}
\label{subsubsec:example1-response}

The LLM opens with generic shoe-selection criteria, embeds the Nike
Pegasus~41 as a \texttt{[Sponsored]} link, and closes with a breathability
tip---consistent with the context guide's instruction to connect product
features to the user's goals.  Once the response is streamed, the orchestrator
fires two asynchronous writes: an impression event appended to the
\texttt{events} table, and a \texttt{daily\_spend\_log} upsert that increments
\texttt{spend} by $\text{bid\_cpm}/1000 = \$0.005$, keeping the budget gate
current for all subsequent requests that day.
